%% NeuroBloom SoftwareX draft based on the original SoftwareX template.
%% This draft currently fills the title, abstract, keywords, and
%% Motivation and significance section while keeping the remaining
%% template structure available for later completion.

\documentclass[preprint,12pt,a4paper]{elsarticle}

\usepackage{amssymb}
\usepackage{algpseudocode}
\usepackage{graphicx}
\usepackage{hyperref}
\usepackage[left=2cm,right=2cm]{geometry}
\setlength{\parindent}{0pt}

\journal{SoftwareX}

\begin{document}



\begin{frontmatter}

\title{NeuroBloom: An open-source platform for longitudinal cognitive task monitoring and clinician-supervised rehabilitation workflows in multiple sclerosis}

\author{Adit Mugdha Das}

\author{Abhijeet Deb Nath}

\author{Kazi Saeed Alam\corref{cor1}}
\ead{saeed.alam@cse.kuet.ac.bd}

\author{Sk. Imran Hossain}

\cortext[cor1]{Corresponding author}

\address{Department of Computer Science and Engineering, Khulna University of Engineering \& Technology, Khulna 9203, Bangladesh}

\begin{abstract}
Multiple sclerosis can affect several cognitive abilities, but gradual changes may be difficult to recognize through occasional clinical assessments. NeuroBloom is an open-source web platform that supports repeated cognitive task sessions and clinician-supervised rehabilitation planning. It integrates 35 tasks across six cognitive domains with pre-session contextual questions on fatigue, sleep, stress, and medication timing. The platform stores trial- and session-level results over time and provides separate interfaces for patients, clinicians, and administrators. Built with Svelte, FastAPI, SQLModel, and PostgreSQL, NeuroBloom supports Bengali and English and can be deployed using Docker Compose. Its analytics module summarizes performance change, response-time variability, longitudinal trends, and associations between contextual factors and task performance. Using synthetic demonstration data, the platform illustrates how repeated cognitive and contextual records can be transformed into descriptive longitudinal indicators and clinician-facing summaries. NeuroBloom provides an open-source and reproducible foundation for longitudinal cognitive monitoring, clinician-supervised rehabilitation workflows, and future digital-biomarker-oriented research in multiple sclerosis.
\end{abstract}


\begin{keyword}
multiple sclerosis \sep digital health \sep cognitive monitoring \sep cognitive rehabilitation \sep digital-biomarker-oriented research \sep clinical software
\end{keyword}

\end{frontmatter}

\begin{table}[!h]
\centering
\begin{tabular}{|l|p{6.5cm}|p{6.5cm}|}
\hline
\textbf{Nr.} & \textbf{Code metadata description} & \textbf{Metadata} \\
\hline

C1 & Current code version & v1.1 \\
\hline

C2 & Permanent link to archived code/repository used for this code version &
Archived version (DOI):
\url{https://doi.org/10.5281/zenodo.21207720}

Source repository:
\url{https://github.com/Adit-Mugdha-das/NeuroBloom} \\
\hline

C3 & Legal Code License & MIT License \\
\hline

C4 & Code versioning system used & git \\
\hline

C5 & Software code languages, tools, and services used &
Python, FastAPI, SQLModel, PostgreSQL, JavaScript, Svelte \\
\hline

C6 & Compilation requirements, operating environments \& dependencies &
Python 3.10+, Node.js 18+, npm, PostgreSQL, modern web browser,
backend dependencies including FastAPI, SQLModel, Uvicorn, and bcrypt,
and frontend build tooling based on SvelteKit/Vite \\
\hline

C7 & If available Link to developer documentation/manual &
\url{https://github.com/Adit-Mugdha-das/NeuroBloom/blob/main/README.md} \\
\hline

C8 & Support email for questions &
saeed.alam@cse.kuet.ac.bd \\
\hline

\end{tabular}
\caption{Code metadata}
\label{codeMetadata}
\end{table}

\section{Motivation and significance}

An estimated 2.9 million people worldwide live with multiple sclerosis
(MS) \cite{atlasms2023}. Although MS is commonly associated with motor
and sensory symptoms, it can also affect processing speed, attention,
working memory, executive function, and cognitive fatigue
\cite{chiaravalloti2008,amato2006,rocca2015}. These cognitive changes
often develop gradually and can be difficult to recognize during
occasional clinical visits. Their severity may also vary with fatigue,
sleep quality, stress, and medication timing, which makes repeated
observation important \cite{amato2006,kalb2018}. Worsening in MS does
not always occur through clearly identifiable relapses. Gradual
progression can occur without an overt clinical attack, including
progression independent of relapse activity (PIRA)
\cite{giovannoni2022,kappos2020}. Subtle changes may therefore remain
unnoticed or poorly documented between formal assessments. Repeated
cognitive-task sessions cannot replace neurological examination or
clinical diagnosis, but they can provide structured records of
performance between clinical visits.

Existing clinical batteries and digital assessment tools provide
valuable approaches for cognitive assessment in MS. In-clinic batteries
such as BICAMS provide validated assessment \cite{benedict2012}, while
remote systems such as Floodlight and Konectom have demonstrated the
feasibility of repeated cognitive testing outside the clinic
\cite{montalban2022,scaramozza2024}. However, repeated cognitive tasks,
adaptive difficulty, contextual self-report, longitudinal storage and
analysis, clinician review, rehabilitation planning, and multilingual
support are not commonly integrated within a single open-source
platform. NeuroBloom was developed to address this software-integration
gap.

The need for accessible follow-up can become even more important in
resource-constrained settings such as Bangladesh. Access to specialist
neurological and cognitive assessment may be limited, while repeated
follow-up can be affected by cost, travel, service availability, and
health literacy \cite{das2017healthliteracy}. Publicly accessible
longitudinal cognitive-monitoring data for Bangladeshi people with MS
also appear to be limited. Language can create an additional barrier,
as many digital cognitive systems are primarily designed for
English-speaking users. These challenges motivated the inclusion of
Bengali alongside English in NeuroBloom. Bengali is among the most
widely spoken languages in the world \cite{ethnologue2025}. Bangladesh
therefore provides an important accessibility context for the platform,
although NeuroBloom itself is designed for broader international use
and future localization to additional languages.
\begin{figure}[!t]
\centering
\includegraphics[width=\textwidth]{figure1}
\caption{Motivation and integrated NeuroBloom workflow for repeated
cognitive-task monitoring and clinician-supervised review.}
\label{fig:motivation_workflow}
\end{figure}


NeuroBloom brings the main parts of this workflow together in one
open-source web platform. Patients can complete repeated cognitive tasks
across six domains with adaptive difficulty, while short pre-session
questions capture contextual factors such as fatigue, sleep, stress,
and medication timing. The platform stores performance over time and
derives descriptive longitudinal indicators, including response-time
variability, fatigue-related patterns, performance trends, and
within-person change. Clinicians can review these records and adjust
rehabilitation plans, while administrators support user management and
platform oversight. These monitoring indicators may also support future
digital-biomarker-oriented research.

Fig.~\ref{fig:motivation_workflow} summarizes the motivation for
NeuroBloom and the integrated workflow connecting repeated cognitive
tasks, contextual data collection, longitudinal monitoring, and
clinician-supervised review. The software architecture and individual
functionalities are described in Section~2.




\section{Software description}

This section describes the design and implementation of NeuroBloom. Section~2.1 presents the role-aware software architecture, including the frontend, backend, analytics, and persistence layers. Section~2.2 describes the main software functionalities, including patient-facing task execution, adaptive session planning, longitudinal analytics, clinician review, and administrative oversight.


\subsection{Software architecture}

NeuroBloom uses a modular, role-aware architecture with four main
layers: the user interface, application services, analytics and
monitoring, and persistent data storage. These layers share the same
backend and data model while supporting separate workflows for
patients, clinicians, and administrators. The overall architecture is
shown in Fig.~\ref{fig:system_architecture}.

\begin{figure}[!t]
\centering
\includegraphics[width=\textwidth]{figure2}
\caption{Role-based system architecture of NeuroBloom. Patient,
clinician, and administrator interfaces share a common FastAPI backend,
analytics layer, and PostgreSQL persistence layer.}
\label{fig:system_architecture}
\end{figure}

The user-interface layer is implemented using Svelte. It provides
separate views for cognitive-task execution, progress tracking,
clinician review, patient management, analytics, and administration.
The patient and clinician interfaces support both Bengali and English
through a custom internationalisation system. Interface text is stored
in a structured translation catalog, while each user's language
preference is saved with the account. When Bengali is selected,
numeric values can be shown using Bengali digits, and selected
letter-based task stimuli are displayed using Bengali script where
appropriate, without changing the underlying task structure or scoring
logic. The current administrative interface is available in English.

The application-services layer is implemented using FastAPI. It manages
authentication, baseline assessment, training sessions, session
planning, reports, messaging, notifications, and coordination between
patients, clinicians, and administrators. A task-planning component also
supports adaptive session generation. It uses baseline performance,
recent task history, and current difficulty levels to identify relevant
domains, reduce immediate task repetition, select suitable tasks, and
update difficulty across later sessions.

The analytics and monitoring layer processes stored task, session, and
contextual records. It derives descriptive measures such as
response-time variability, fatigue-related patterns, longitudinal
performance trends, and within-person change. These outputs are
presented for monitoring and clinician review rather than being used as
an independent diagnostic system. The processing flow from task
interaction and contextual information to longitudinal indicators is
shown in Fig.~\ref{fig:biomarker_pipeline}.

\begin{figure}[!ht]
\centering
\includegraphics[width=\textwidth]{figure3}
\caption{Processing pipeline from cognitive task interaction and
contextual self-report to descriptive longitudinal indicators and
clinician review.}
\label{fig:biomarker_pipeline}
\end{figure}

The data layer uses SQLModel and PostgreSQL to store user information,
task results, session history, and contextual records such as fatigue,
sleep quality, stress, and medication timing. Task records can include
score, accuracy, response time, completion information, and difficulty
level. Keeping these records in a shared persistent data model allows
patient performance to be reviewed across multiple sessions rather than
as isolated task results.

The separation of interface, application logic, analytics, and storage
makes the software easier to maintain and extend as new task modules,
analytics features, or language support are added. NeuroBloom also
supports reproducible deployment through Docker Compose. Backend and
frontend dependencies are documented in the repository, and GitHub
Actions runs the configured test suites on each commit. The repository
includes numerical scoring fixtures for all 35 tasks, task-inventory
checks, backend API tests, and frontend unit tests. The archived software release (v1.1) is
publicly available through Zenodo.


\subsection{Software functionalities}

\subsubsection{Patient interaction and session execution}

NeuroBloom provides a patient-facing workflow that begins with an initial
baseline assessment and continues through repeated training sessions. On
first use, the patient completes a baseline assessment across all six
cognitive domains, and the resulting scores are used to set initial task
difficulty levels and identify areas for greater training focus. The patient
can then search the clinician directory and send an assignment request; once
accepted, the clinician can review the patient's baseline results and later
session history. Patients subsequently complete structured training tasks,
review their session progress, and optionally provide pre-session contextual
information such as fatigue, sleep quality, stress, and medication timing.
The patient dashboard also shows training history, streaks, achievements,
and badges to support continued engagement.
The patient interface supports both Bengali and English. Representative
cognitive-task and result interfaces in both languages are presented in the
Cognitive task framework subsection.

\begin{figure}[!t]
\centering

\begin{minipage}[t]{0.48\textwidth}
\centering
\includegraphics[width=\textwidth]{fig4part1}\\[-0.2em]
{\small (a)}
\end{minipage}
\hfill
\begin{minipage}[t]{0.48\textwidth}
\centering
\includegraphics[width=\textwidth]{fig4part2}\\[-0.2em]
{\small (b)}
\end{minipage}

\vspace{0.4em}

\begin{minipage}[t]{0.48\textwidth}
\centering
\includegraphics[width=\textwidth]{fig4part3}\\[-0.2em]
{\small (c)}
\end{minipage}
\hfill
\begin{minipage}[t]{0.48\textwidth}
\centering
\includegraphics[width=\textwidth]{fig4part4}\\[-0.2em]
{\small (d)}
\end{minipage}

\caption{Examples of cognitive tasks in the English NeuroBloom interface.
(a) Interaction screen for the first cognitive task.
(b) Corresponding result screen for the first task.
(c) Interaction screen for the second cognitive task.
(d) Corresponding result screen for the second task.}
\label{fig:english_tasks}
\end{figure}

\subsubsection{Cognitive task framework}

NeuroBloom includes 35 cognitive tasks across six domains: working memory,
processing speed, attention, cognitive flexibility, planning, and visual
scanning. The tasks use a common session model that records score, accuracy,
response time, completion behavior, and difficulty progression. Before each
cognitive task begins, NeuroBloom presents an instruction screen describing
the task objective and required interaction. Selected tasks also include a
short practice round before the scored task begins, allowing users to become
familiar with the task procedure without including the practice responses in
the recorded task result. Each task supports 10 difficulty levels, allowing
patient performance to be followed across repeated sessions. Examples of two cognitive tasks and their corresponding result screens are
shown in English in Fig.~\ref{fig:english_tasks} and in Bengali in
Fig.~\ref{fig:bengali_tasks}. The Bengali interfaces preserve the underlying
task structure and scoring logic while localizing the user-facing
presentation.


\begin{figure}[!t]
\centering

\begin{minipage}[t]{0.48\textwidth}
\centering
\includegraphics[width=\textwidth]{fig5part1}\\[-0.2em]
{\small (a)}
\end{minipage}
\hfill
\begin{minipage}[t]{0.48\textwidth}
\centering
\includegraphics[width=\textwidth]{fig5part2}\\[-0.2em]
{\small (b)}
\end{minipage}

\vspace{0.4em}

\begin{minipage}[t]{0.48\textwidth}
\centering
\includegraphics[width=\textwidth]{fig5part3}\\[-0.2em]
{\small (c)}
\end{minipage}
\hfill
\begin{minipage}[t]{0.48\textwidth}
\centering
\includegraphics[width=\textwidth]{fig5part4}\\[-0.2em]
{\small (d)}
\end{minipage}

\caption{Bengali versions of the cognitive tasks shown in
Fig.~\ref{fig:english_tasks}.
(a) Bengali interaction screen for the first cognitive task.
(b) Corresponding result screen for the first task.
(c) Bengali interaction screen for the second cognitive task.
(d) Corresponding result screen for the second task.}
\label{fig:bengali_tasks}
\end{figure}

Before repeated training begins, the baseline assessment establishes
domain-level performance scores. These scores are used to initialize task
difficulty and identify areas for greater training focus. During later
sessions, the planning mechanism uses baseline performance, recent task
history, and current difficulty levels to prioritize weaker domains, select
tasks that have not been used recently where possible, and update difficulty
after completed sessions. Four tasks are selected for each planned session.
The overall process is summarized in Algorithm~1.

\begin{center}
\begin{minipage}{0.92\linewidth}
\small
\textbf{Algorithm 1. Personalized session planning and task rotation in NeuroBloom}
\begin{algorithmic}[1]
\Require Baseline scores $B[d]$, recent task history $H[d]$, current difficulties $L[d]$
\State $\textit{sorted\_domains} \gets \Call{SortAscending}{B}$
\State $\textit{primary\_focus} \gets \textit{sorted\_domains}[0:2]$
\State $\textit{secondary\_focus} \gets \textit{sorted\_domains}[2:4]$
\State $\textit{maintenance} \gets \textit{sorted\_domains}[4:6]$
\ForAll{domain $d$ in $B$}
    \State $L[d] \gets \Call{ScoreToDifficulty}{B[d]}$
\EndFor
\State $\textit{session\_domains} \gets
\textit{primary\_focus} \cup \textit{secondary\_focus}$
\ForAll{domain $d$ in $\textit{session\_domains}$}
    \State $\textit{available} \gets \Call{GetAvailableTasks}{d}$
    \State $\textit{recent} \gets \Call{GetRecentTasks}{H[d], d}$
    \State $\textit{candidates} \gets
    \Call{ExcludeRecentWhenPossible}{\textit{available}, \textit{recent}}$
    \State $\textit{strategy} \gets \Call{RotationStrategy}{d}$
    \State $\textit{task}[d] \gets
    \Call{SelectTask}{\textit{candidates}, \textit{strategy}, L[d]}$
\EndFor
\ForAll{completed task result $r$}
    \State \Call{StoreTrainingSession}{$r$}
    \State $L[r.domain] \gets
    \Call{UpdateDifficulty}{L[r.domain], r.score, r.accuracy, r.reaction\_time}$
\EndFor
\If{\Call{ShouldRebalanceFocusAreas}{active training plan}}
    \State \Call{RebalanceFocusAreas}{active training plan}
\EndIf
\State \Return personalized four-task session plan
\end{algorithmic}
\end{minipage}
\end{center}

\subsubsection{Descriptive longitudinal indicators for
digital-biomarker-oriented research}

Beyond individual task results, NeuroBloom derives descriptive indicators
from accumulated trial, session, and contextual records. The current
analytics layer uses deterministic statistical calculations based on
trial-level accuracy and reaction time, session-level scores, and contextual
self-report variables. Four main constructs are produced: a fatigue-related
proxy, reaction-time variability, longitudinal trend, and within-person
change. Contextual associations are also calculated. These outputs are used
for longitudinal monitoring and clinician review and may support future
digital-biomarker-oriented research.

\textbf{Fatigue-related proxy.}
Within-session change is summarized by comparing performance during the
first and last quartiles of valid trials:

\begin{equation}
\Delta Acc = Acc_{Q1} - Acc_{Q4}, \qquad
\Delta RT = RT_{Q4} - RT_{Q1},
\end{equation}

where $Q1$ and $Q4$ denote the first and last quartiles of trials in a
session. A trial-wise linear slope $\beta$ of accuracy over time is also
computed. The normalized proxy is defined as

\begin{equation}
F = 0.5A + 0.3R + 0.2S,
\end{equation}

where
$A=\min(1,\max(0,\Delta Acc/30))$,
$R=\min(1,\max(0,\Delta RT/200))$, and
$S=\min(1,\max(0,-100\beta))$.
This is a configurable engineering proxy rather than a clinically calibrated
fatigue measure. The default weights give greater contribution to accuracy,
while the normalization constants are developer-selected saturation values.
Sessions with insufficient valid trial data return insufficient data rather
than a fatigue-proxy value.

\textbf{Reaction-time variability.}
NeuroBloom summarizes intra-individual variability in reaction time using
the coefficient of variation:

\begin{equation}
CV = \frac{\sigma_{RT}}{\mu_{RT}},
\end{equation}

where $\sigma_{RT}$ and $\mu_{RT}$ are the standard deviation and mean of
valid reaction times. Complementary measures, including median absolute
deviation (MAD), interquartile range (IQR), and full range, are also
calculated to describe the spread of reaction times without relying on a
single variability statistic.

\textbf{Longitudinal trend.}
Repeated performance values are summarized using an exponentially weighted
moving average (EWMA):

\begin{equation}
EWMA_t = \alpha x_t + (1-\alpha)EWMA_{t-1},
\end{equation}

with $\alpha=0.2$ in the current implementation. Trend direction and
strength are estimated from the slope of the five most recent EWMA points.
The resulting direction is shown in the clinician-facing interface as
improving, stable, or declining.

\textbf{Within-person standardized change.}
Because independent test--retest reliability estimates are not currently
available for the implemented task versions, NeuroBloom reports a
within-person standardized change score (WPSC) rather than a reliable
change index:

\begin{equation}
WPSC =
\frac{X_{\mathrm{latest}} - X_{\mathrm{baseline}}}
{\sqrt{2}\,s_{\mathrm{within}}},
\end{equation}

where $X_{\mathrm{baseline}}$ and $X_{\mathrm{latest}}$ are the first and
latest valid comparable scores, and $s_{\mathrm{within}}$ is the
within-person sample standard deviation. At least three chronologically
ordered scores with non-zero variance are required. Missing and non-finite
values are excluded, while outliers are retained. If these requirements are
not satisfied, the interface reports insufficient data. WPSC is descriptive
and has no threshold for clinical significance.

\textbf{Contextual associations.}
The platform also explores relationships between performance and
pre-session information, including fatigue level, sleep quality, time since
medication, stress level, pain level, and readiness level. Pearson
correlation coefficients are calculated with session score and mean reaction
time and are displayed as exploratory associations rather than causal or
clinical relationships.  Together, these indicators provide statistically
derived monitoring signals that help clinicians interpret longitudinal performance changes and
adjust rehabilitation plans over time.

\subsubsection{Clinician-facing review and rehabilitation planning}

The clinician interface supports patient review, rehabilitation planning,
and communication. Once a patient has connected with a clinician, the
clinician can review baseline results, recent task performance, adherence,
longitudinal indicators, trends, and progress reports. Based on this
information, the clinician can modify the patient's rehabilitation plan by
changing task focus, difficulty levels, and planned session frequency.The platform can also generate configurable review flags when predefined
changes in performance or adherence are detected. These flags draw attention
to records that may require review and are visible to both clinicians and
administrators. A bidirectional messaging system allows patients to report concerns between
sessions and clinicians to respond through the platform. Patients and
clinicians can also download performance records and progress reports for
review, documentation, or follow-up discussion. Together, these functions
connect repeated performance monitoring, communication, and
clinician-supervised rehabilitation planning within a single workflow. The
clinician patient-review and longitudinal rehabilitation-planning
interfaces are shown in Fig.~\ref{fig:clinical_admin_workflow}(a)--(b).


\subsubsection{Administrative oversight and governance}

NeuroBloom provides administrative functions for managing the platform in
a multi-user environment. Administrators can review system-level activity
and system-health information, manage users and role relationships, oversee
notification delivery, inspect audit logs for accountability and
traceability, and export anonymized research-oriented data for downstream
analysis. These functions provide platform-level oversight while keeping
patient-performance interpretation and rehabilitation planning within the
clinician workflow. The analytics and administrative interfaces are shown in
Fig.~\ref{fig:clinical_admin_workflow}(c)--(d).

\begin{figure}[!t]
\centering

\begin{minipage}[t]{0.48\textwidth}
\centering
\includegraphics[width=\textwidth]{fig6part1}\\[-0.2em]
{\small (a)}
\end{minipage}
\hfill
\begin{minipage}[t]{0.48\textwidth}
\centering
\includegraphics[width=\textwidth]{fig6part2}\\[-0.2em]
{\small (b)}
\end{minipage}

\vspace{0.4em}

\begin{minipage}[t]{0.48\textwidth}
\centering
\includegraphics[width=\textwidth]{fig6part3}\\[-0.2em]
{\small (c)}
\end{minipage}
\hfill
\begin{minipage}[t]{0.48\textwidth}
\centering
\includegraphics[width=\textwidth]{fig6part4}\\[-0.2em]
{\small (d)}
\end{minipage}

\caption{Clinician, analytics, and administrative interfaces in NeuroBloom.
(a) Clinician patient-list and review interface.
(b) Longitudinal monitoring and rehabilitation-planning view.
(c) Longitudinal monitoring and contextual analytics dashboard.
(d) Administrative interface for system monitoring and governance.}
\label{fig:clinical_admin_workflow}
\end{figure}

\section{Illustrative examples}

To demonstrate the complete NeuroBloom workflow, we provide a synthetic
patient example using demonstration data only. The patient first completes
the six-domain baseline assessment, and the resulting scores are used to set
initial difficulty levels and identify areas for greater training focus. The
patient then searches the clinician directory and sends an assignment
request. After the clinician accepts the request, the clinician can review
the baseline results and subsequent session history. Based on the baseline
scores, recent task history, and current difficulty levels, NeuroBloom
selects four tasks for the next session. Before starting, the patient also
records contextual information such as fatigue, sleep quality, stress, and
medication timing.

During the session, NeuroBloom stores task-level score, accuracy, response
time, completion information, and difficulty level. These records are
processed by the analytics module to generate descriptive longitudinal
indicators. For example, response-time variability can be calculated from
valid reaction times using the coefficient of variation described in
Section~2.2.3. The clinician can then review the patient's task results,
contextual information, longitudinal indicators, and progress history. When
appropriate, the clinician can modify the rehabilitation plan by adjusting
task focus, difficulty, or planned session frequency. Administrative users
support the workflow through user and role management, system monitoring,
audit review, and research-oriented data export.

All values used in this example are synthetic and are intended only to
demonstrate software functionality. They do not represent a real patient or
clinical outcome. The source code and reproducibility materials are
available in the NeuroBloom repository at
\url{https://github.com/Adit-Mugdha-das/NeuroBloom}, and the corresponding
archived software release is available at
\url{https://doi.org/10.5281/zenodo.21207720}. A supplementary video further
demonstrates the patient, clinician, and administrative workflows and
complements the illustrative example.

\section{Impact}


NeuroBloom has potential impact beyond longitudinal data collection because
it provides a structured environment for repeated cognitive activity and
clinician-supervised rehabilitation. Its 35 tasks are original digital
implementations inspired by established cognitive paradigms, including
N-back, digit span, serial addition, Stroop-type interference, Go/No-Go,
trail-making, card sorting, planning, and visual-search paradigms. Together,
these tasks engage cognitive functions that are commonly affected in MS,
including working memory, processing speed, attention, executive control,
planning, and visual attention. Previous randomized studies of computerized
cognitive rehabilitation in people with MS have reported improvements in
selected cognitive functions, including processing speed, attention,
working memory, and executive performance
\cite{hancock2015,chiaravalloti2022,naeeni2022}. NeuroBloom does not assume
that its individual task implementations are equivalent to standardized
clinical instruments; rather, it provides an accessible software framework
through which established cognitive-training principles can be delivered
repeatedly, adapted over time, and studied under clinician supervision.

A practical impact of this approach is the possibility of extending
structured cognitive activity beyond occasional face-to-face assessment.
Patients can continue planned cognitive sessions between clinical visits,
while clinicians can review performance history, adherence, contextual
information, and longitudinal changes when deciding how to adjust subsequent
rehabilitation activity. This creates a closer connection between repeated
home-based cognitive activity and professional review than would be possible
when training and follow-up are performed as separate processes. It may be
particularly useful for patients for whom frequent specialist visits are
difficult, while preserving the clinician's role in interpreting the
information and supervising rehabilitation decisions.

The accessibility impact is particularly relevant to resource-constrained
and multilingual settings. Commercial or institution-specific digital
rehabilitation systems may be difficult to adopt where financial,
infrastructural, language, or specialist-access barriers exist. NeuroBloom
is open source, browser based, and provides Bengali as well as English
interfaces, allowing the same software to be deployed without requiring a
proprietary platform. In Bangladesh, where health-literacy and healthcare
access challenges have been documented \cite{das2017healthliteracy}, this
design can provide a practical foundation for locally accessible cognitive
rehabilitation research and supervised follow-up. Bengali support is also
relevant beyond Bangladesh because Bengali is among the world's most widely
spoken languages \cite{ethnologue2025}. The localization framework can be
extended to other languages, making the underlying approach applicable to
other underserved populations.

NeuroBloom may also have an important research impact. Repeated task
performance, contextual information, difficulty progression, adherence, and
clinician-supervised rehabilitation history can be retained within the same
longitudinal record. This creates opportunities to investigate how cognitive
performance changes over time, how different individuals respond to repeated
training, and how factors such as fatigue, sleep, stress, and medication
timing are associated with performance. Such data can support future
digital-biomarker-oriented research \cite{insel2017,coravos2019} and may
also contribute to the development of locally relevant longitudinal
datasets in populations that are currently underrepresented in digital
health research.

Finally, the open-source design makes the impact of NeuroBloom extend beyond
its current MS-focused implementation. Researchers can introduce new task
paradigms, rehabilitation strategies, analytical methods, languages, and
disease-specific workflows while retaining the same patient--clinician
infrastructure. The platform can therefore serve both as an operational
cognitive-rehabilitation software system and as an experimental foundation
for studying new approaches to longitudinal cognitive support. Future pilot
and clinical studies can evaluate usability, adherence, engagement,
cognitive outcomes, and the effectiveness of the complete
clinician-supervised workflow in real-world settings.


\section{Conclusions}

NeuroBloom addresses a software-integration gap in MS cognitive monitoring
and rehabilitation planning by combining repeated cognitive task execution,
contextual data capture, adaptive session planning, longitudinal analytics,
and clinician-facing review within a single open-source platform. Its
modular architecture and Bengali--English interface provide a reproducible
foundation for digital health research and supervised follow-up, including
use in resource-constrained settings.Future work will focus on prospective clinical evaluation, assessment of longitudinal engagement and real-world usability, expansion of the cognitive task library, and further development of the monitoring and analytics
components. Clinical studies will also be needed to evaluate practice
effects, test--retest reliability, and the influence of device and
environmental conditions on repeated task performance. NeuroBloom is
currently a research software platform and is not intended to provide
diagnosis, replace clinical judgment, or function as a validated medical
device.

\section*{CRediT authorship contribution statement}

\textbf{Adit Mugdha Das:} Software, Data curation, Methodology, Validation, Visualization, Writing -- original draft. \textbf{Abhijeet Deb Nath:} Methodology, Software, Data curation. \textbf{Kazi Saeed Alam:} Conceptualization, Supervision, Resources, Validation, Writing--review and editing. \textbf{Sk. Imran Hossain:} Supervision, Writing--review and editing.

\section*{Data availability}

No real patient data were collected or analyzed in this study. The
manuscript uses only synthetic and demonstration data to illustrate the
software workflow. The source code for NeuroBloom is publicly available at
\url{https://github.com/Adit-Mugdha-das/NeuroBloom}, and the archived
software release used in this manuscript is available through Zenodo at
\url{https://doi.org/10.5281/zenodo.21207720}.

\section*{Acknowledgements}
The authors would like to acknowledge the support of the Department of Computer Science and Engineering, Khulna University of Engineering \& Technology, Khulna, Bangladesh, for providing the necessary resources and infrastructure to carry out this project.

\section*{Ethics and data privacy statement}

No real patient data were collected, analyzed, or reported in this manuscript. All screenshots and examples use synthetic or demonstration data; therefore, ethical approval and informed consent were not required for this software description. Any future clinical deployment of NeuroBloom will require appropriate ethical approval, informed consent, and privacy-preserving handling of patient data.

\section*{Declaration of Generative AI and AI-assisted technologies in the writing process}
During the preparation of this work the authors used AI tools in order to improve the readability and language of the manuscript. After using this service, the authors reviewed and edited the content as needed and take full responsibility for the content of the published article.

\section*{Declaration of competing interest}
The authors declare that they have no known competing financial interests or personal relationships that could have appeared to influence the work reported in this paper.

\begin{thebibliography}{00}

\bibitem{atlasms2023}
Multiple Sclerosis International Federation,
Atlas of MS: Number of people with MS, 2023.
\url{https://atlasofms.org/map/global/epidemiology/number-of-people-with-ms}
(accessed 7 August 2026).

\bibitem{chiaravalloti2008}
N.D.~Chiaravalloti, J.~DeLuca,
Cognitive impairment in multiple sclerosis,
Lancet Neurol. 7 (2008) 1139--1151.

\bibitem{amato2006}
M.P.~Amato, V.~Zipoli, E.~Portaccio,
Multiple sclerosis-related cognitive changes: a review of cross-sectional
and longitudinal studies,
J. Neurol. Sci. 245 (2006) 41--46.

\bibitem{rocca2015}
M.A.~Rocca, M.P.~Amato, N.~De Stefano, et al.,
Clinical and imaging assessment of cognitive dysfunction in multiple
sclerosis,
Lancet Neurol. 14 (2015) 302--317.

\bibitem{kalb2018}
R.~Kalb, M.~Beier, R.H.B.~Benedict, et al.,
Recommendations for cognitive screening and management in multiple
sclerosis care,
Mult. Scler. 24 (2018) 1665--1680.

\bibitem{giovannoni2022}
G.~Giovannoni, V.~Popescu, J.~Wuerfel, et al.,
Smouldering multiple sclerosis: the `real MS',
Ther. Adv. Neurol. Disord. 15 (2022) 17562864211066751.

\bibitem{kappos2020}
L.~Kappos, J.S.~Wolinsky, G.~Giovannoni, et al.,
Contribution of relapse-independent progression vs relapse-associated
worsening to overall confirmed disability accumulation in typical relapsing
multiple sclerosis in a pooled analysis of 2 randomized clinical trials,
JAMA Neurol. 77 (2020) 1132--1140.

\bibitem{benedict2012}
R.H.B.~Benedict, M.P.~Amato, J.~Boringa, et al.,
Brief International Cognitive Assessment for Multiple Sclerosis (BICAMS):
international standards for validation,
BMC Neurol. 12 (2012) 55.

\bibitem{montalban2022}
X.~Montalban, J.~Graves, L.~Midaglia, et al.,
A smartphone sensor-based digital outcome assessment of multiple sclerosis,
Mult. Scler. 28 (2022) 654--664.

\bibitem{scaramozza2024}
M.~Scaramozza, P.A.~Chiesa, L.~Zajac, et al.,
Konectom cognitive processing speed test enables reliable remote,
unsupervised cognitive assessment in people with multiple sclerosis:
exploring the use of substitution time as a novel digital outcome measure,
Mult. Scler. 30 (2024) 1193--1204.

\bibitem{das2017healthliteracy}
S.~Das, M.N.~Mia, S.M.A.~Hanifi, S.~Hoque, A.~Bhuiya,
Health literacy in a community with low levels of education: findings from
Chakaria, a rural area of Bangladesh,
BMC Public Health 17 (2017) 203.
\url{https://doi.org/10.1186/s12889-017-4097-y}.

\bibitem{ethnologue2025}
D.M.~Eberhard, G.F.~Simons, C.D.~Fennig (Eds.),
Bengali, in: Ethnologue: Languages of the World, 28th ed.,
SIL International, Dallas, Texas, 2025.
\url{https://www.ethnologue.com/language/ben/}
(accessed 7 August 2026).



\bibitem{insel2017}
T.R.~Insel,
Digital phenotyping: technology for a new science of behavior,
JAMA 318 (2017) 1215--1216.

\bibitem{coravos2019}
A.~Coravos, S.~Khozin, K.D.~Mandl,
Developing and adopting safe and effective digital biomarkers to improve
patient outcomes,
npj Digit. Med. 2 (2019) 14.

\bibitem{tseriotis2026}
V.-S.~Tseriotis, L.M.~Peeters, L.~Geys, et al.,
Digital patient experience tools in multiple sclerosis: a landscape analysis
of the global Patient-Reported Outcomes in Multiple Sclerosis (PROMS)
initiative,
eClinicalMedicine 94 (2026) 103821.

\bibitem{hancock2015}
L.M.~Hancock, J.M.~Bruce, A.S.~Bruce, S.G.~Lynch,
Processing speed and working memory training in multiple sclerosis:
a double-blind randomized controlled pilot study,
J. Clin. Exp. Neuropsychol. 37 (2015) 113--127.
\url{https://doi.org/10.1080/13803395.2014.989818}.

\bibitem{chiaravalloti2022}
N.D.~Chiaravalloti, S.L.~Costa, N.B.~Moore, K.~Costanza, J.~DeLuca,
The efficacy of speed of processing training for improving processing
speed in individuals with multiple sclerosis: a randomized clinical trial,
J. Neurol. 269 (2022) 3614--3624.
\url{https://doi.org/10.1007/s00415-022-10980-9}.

\bibitem{naeeni2022}
M.N.~Davarani, A.A.~Darestani, P.~Hassani-Abharian, et al.,
RehaCom rehabilitation training improves a wide-range of cognitive
functions in multiple sclerosis patients,
Appl. Neuropsychol. Adult 29 (2022) 262--272.
\url{https://doi.org/10.1080/23279095.2020.1747070}.

\end{thebibliography}

\end{document}
